\section{Equivalence at Matched Dimension}
\label{sec:equivalence}

The pair $S_4$/$A_5$ was designed into the family as the decisive
contrast: the same $\dmin = 3$, opposite sides of the solvability
divide, the axis on which expressivity theory separates state-tracking
architectures \citep{merrill2024illusion, grazzi2025negative,
siems2025deltaproduct}. If learned representational dimension were set
by complexity class, the pair should separate; if by representation
dimension, it should coincide. Because the interesting outcome is a
null, the comparison was pre-registered as an equivalence test rather
than a difference test: Welch two-one-sided tests on restricted
effective rank at margin $\pm 0.5$ rank-units (half the spacing of the
$\dmin$ ladder), $n = 5$ seeds per group, with a pre-run power
simulation showing 100\% power to \emph{reject} equivalence at a true
gap of 1.0 rank-units, so a real class effect of one ladder step could
not have hidden inside the margin.

The observed difference is $+0.019$ rank-units (se 0.037, df 7.8), and
both one-sided tests pass at roughly seven times the critical value
($t = 13.06$ and $14.12$ against $t_{\mathrm{crit}} = 1.865$):
equivalence is declared decisively.
% <!-- evidence: U2 -->
The two groups' seed clusters are visually coincident in
Figure~\ref{fig:convergence} (inset). Convergence here is governed by
what the algebra demands, not by how hard the class is to compute.
